\chapter{Case study 2 - quantum physics}

The second case study of a particular model of physics is indeed, quantum physics. Despite such success in modelling and analysing the world, classical physics hits deadlock where behaviours and observations are so little, effects on macroscopic scale are so large, and in some case the system is too small to isolate such that results and predictions come in fruitless and of particular failure. Of such, solution to such problem comes in the form of creating a new theory, connected still in the sense of certain special cases or edge cases, and of removing factors, with \textit{new postulate} and structure accompanied, and of models that would be fairly difficult to incorporate, aside from of modern classical physics. Such is similar to how relativity is done, just with a generalization of missing content. 

The mathematical technique for modelling quantum physics is also fairly different from classical physics, mainly in the way that the quantum mechanical system is treated and how is that formed. Specifically, one will expect to master of linear algebra, mathematical analysis, topology - especially Hilbert space and rigged Hilbert space topology, differential geometry (kinda), manifold analysis, operator theory, tensor theory, probability theory, and group theory, especially Galois group. Such is not even enough for certain niche covering of quantum physics, but that is for the edge case interpretation they require to express such system of interest (not like I want to discuss of field theory right from the start). For said motion, we can concur on an observation that quantum system behaves toward a much, very much different axioms set, and said thing is what we need to cover and analyse of our own interest. 

\section{Mathematical preliminaries}

There are certainly some specification of knowledge when it comes to understanding quantum mechanics. For now, those mathematical concepts might look familiar to somewhere else, but they contain little to no physics of such. Toward that end, let us take an inquiry on whatever the heck is required to understand this eldritch horror.
\section{Linear and inner vector spaces}

By itself, a \textit{vector space} is a linear space. That is why it is even the main structure of the linear algebra course. Of course, again, there exist the nonlinear manifold, the \textit{affine space} and so on, but they are special case of dropping certain linearity criteria. We would then review such concept here for reference purpose. 

\begin{definition}[(Linear) vector space]
    A vector space $\mathbb{V}$ is a set of objects called \textit{vectors}, denoted $v,\vec{v},\mathbf{v}$ or $\ket{v}$, over a field $\mathbf{F}$, equipped with two operators $(+,\cdot)$ for addition and scalar multiplication, such that the following holds: 
    \begin{enumerate}[topsep=1pt,itemsep=1pt,leftmargin=1cm]
        \item Closure: if $V,W\in \mathbb{V}$, then $V+W\in\mathbb{V}$. 
        \item Closure of scalar multiplication: if $V\in \mathbb{V},\lambda\in \mathbb{F}$, then $\lambda V\in \mathbb{V}$. 
        \item Addition commutativity: if $V,W\in \mathbb{V}$, then $V+W=W+V$. 
        \item Addition associativity: if $V,W\in \mathbb{V}$, then $V+W+Z=(V+W)+Z$.
        \item Scalar multiplication associativity: if $a\in \mathbb{F}$ and $V\in\mathbb{V}$, then $a(bV)=(ab)V$. 
        \item Scalar multiplication distributivity: if $a,b\in\mathbb{F}$, and $V,W\in\mathbf{V}$, then $(a+b)V=aV+bV$ and $a(V+W)=aV+aW$. 
        \item Existence of nullity: There exists a \textit{null vector}, denoted $0,\vec{0},\mathbf{0}$ or $\ket{0}$ such that $V+0=V$ and $0V =0$, for $V\in\mathbb{V}$. 
        \item Existence of negatives: There exists for every $v\in\mathbb{V}$ a vector, written as $-v$, such that $v+(-v)=0$. 
        \item Unitary: There exists the identity vector, denoted $1,\mathbf{1},\ket{1}$, such that $1v=v$ for all $v\in\mathbb{V}$. 
    \end{enumerate}
\end{definition}
Here, we make use of a new notation, which stems from the `official' Dirac braket notation, $\ket{v}$ to denote such vector with name. Think of it as an abstract object in such space, and we now only write it like that so it would be familiar later on. for any mathematical structure that is embedded and structured to take form of a vector space, then each individual vector represents certain subject of said structure, for example, the \textit{state} of the structure in particular configuration. If the field itself is replaced of the real number $\mathbb{R}$, then we say it as a real vector space. Otherwise, if the field is the complex field, $\mathbb{C}$, then it is said to be a complex vector space. This topic should be treated on its own section rather than leaving it only here. 

Such encoding allows for two types of general interest vector space. One is the \textit{discrete basis coordinate space}, for general vectors like $v=(x_1,x_2,x_3,\dots)$, or it can be the space of certain \textit{functions} of some type, for example the space of all differentiable functions. In such case, the field still holds as scalar, but now, we have the entire function as a particular basis or coordinate. Usually, in a very mathy term, it is called a \textbf{function space}, and it is a topological space whose points are now functions instead, more like a wrapper above the typical structure. 

Next, our definition of vector space consequently leave us with a construction of a \textit{basis}. Basically, this mean that they possess bases such that for a vector space $E$, there exist some (one or more) set of vectors $\{e_{1},e_{1},\dots,e_{n}\}$ of the same dimension, such that every vector $\ket{v}\in E$ can be represented by 
\begin{equation}
    \ket{v} = \lambda_{1}e_1 + \lambda_{2}e_{2} + \dots + \lambda_{n} e_{n}
\end{equation}
for all $(\lambda_{1},\dots,\lambda_{n}\in \mathbb{R})$. This is defined as the \textit{linear combination} of such basis upon to express the vector $\ket{v}$. We note that this works for, in general, the basis of real space. The same happens with complex space, but is much, much more difficult to analyze. This definition of a basis, on such is valid for \textit{finite-dimensional space}. We will see how we can extend such notion to the case of infinitely dimensional space later on. For a set of vectors, now we might as well denote here of this physics section as $\ket{\cdot}$, so $\ket{1},\ket{2},\dots,\ket{n}$, is said to be used to construct a \textit{linear combination} on vector space $E$ if , denoting $(a_{i})_{i\in I}$ for index set $I$ the \textit{$I$-indexed family} of element of a set $A$ to be the function $a:I\to A$, the \textit{subfamily} of $(u_{i})_{i\in I}$ being the family $(u_{j})_{j\in J}$ for any $J\subset I$, satisfying the following: 
\begin{definition}[Linear combination]
    Let $E$ be a vector space. A vector $v\in E$ is a \textit{linear combination of a family} $(u_{i})_{i\in I}$ of elements of $E$ iff there is a family $(\lambda_{i})_{i\in I}$ of scalars in $\mathbb{R}$ such that 
    \begin{equation}
        v = \sum_{i\in I} \lambda_{i} u_{i}
    \end{equation}
    When $I=\varnothing$, we stipulate that $v=0$. 
\end{definition}
Then, from such definition, we can then define \textit{linear (in)dependence} of such vector set. 
\begin{definition}[Linearly (in)dependent]
    We say that a family $(u_{i})_{i\in I}$ is \textit{linearly independent} iff for every family $(\lambda_{i})_{i\in I}$ of scalar in $\mathbb{R}$, 
    \begin{equation}
        \sum_{i\in I} \lambda_{i}u_{i} = 0 \Rightarrow \lambda_{i} \ne 0 \: \forall i\in I 
    \end{equation}
    That is, the linear set is unique such that only the null vector yields such correspondence. Equivalently, a family $(u_{i})_{i\in I}$ is said to be \textit{linearly dependent} iff there is some family $(\lambda_{i})_{i\in I}$ of scalars in $\mathbb{R}$ such that 
    \begin{equation}
        \sum_{i\in I} \lambda_{i} u_{i} = 0 \Rightarrow \lambda_{j} \neq 0 , j \in I
    \end{equation}
    We agree also that for $I = \varnothing$, then the family $\varnothing$ is linearly independent. 
\end{definition}

The above simply state that it is linear independent iff either $I\neq \varnothing$, or $I$ consists of a single element $i$ and $u_i \neq 0$, or $|I|\leq 2$ and no vector $u_{j}$ in the family can be expressed as a linear combination of other vectors in the family. The reverse happen for linear dependent. Thus, we can see that, for something to be a basis, it needs to be linearly independent. However, we also look at the criteria that any vector $v\in E$ should be able to be expressed by such linear combination. Then, the set of linear independent vectors must also \textit{span} the vector space $E$, or in simpler term, of the same dimension as vector space $E$ indict upon. We then also want to express now the late definition of what constitute the \textit{dimension} of such space. 

\begin{definition}[Dimension]
    A vector space has \textit{dimension} $n$ if it can accommodate a maximum of $n$ linearly independent vectors. For a vector space $V$, this is denoted by $V^{n}(\mathbb{F})$ follow by their field, which is not necessarily real or complex. 
\end{definition}

Moving on with vector space, without the usual vector subspaces and stuff (which has been covered in its respective chapter on linear algebra), we can add additional structure to the space itself. One of such is the \textit{inner product}, or the scalar product for a linear vector space. The inner product is oddly similar to the dot product in the usual expression of vector algebra, such is $\langle A , B \rangle = |A||B|\cos{(\theta)}$. This expression, in coordinate has the form of $\langle A , B \rangle = A_{x}B_{x}+A_{y}B_{y}$ and else, which can be generalized to $n$-space. The formula then indeed specify that dot product as a functional $\langle \cdot, \cdot \rangle : V\times V \to \mathbb{F}$ of the embedding it receives. To this end, it satisfies the symmetry property, i.e. $\langle A , B \rangle = \langle B,A \rangle $, positive semidefiniteness\footnote{The term here might be confusing, but this might help. For a specific function, then, for now is the inner product, being positive definite means that $\langle A , B \rangle >0$ whenever $A\neq 0$. Positive \textit{semidefinite} means that now $\langle A ,A \rangle \geq 0$ instead for all $A$, and could be zero even if $A\neq 0$. In matrix notation, this means you indict that $x^{\top}Ax\geq 0$ for all $x$, and that $A$ symmetric matrix is what called an $A$ positive semidefinite matrix. The only reason that it is there, is because the inner product will ultimately encode the more general notion of dot product, which makes it introducing new geometry on top. Simple case, $A=I$ of identity matrix on $n$-space.}, which is $\langle A , A \rangle \geq 0$, being $0$ iff $A=0$ for $A\in E$, and linearity property 
\begin{equation*}
    \langle A , (\lambda B + \psi C) \rangle = \lambda \langle A , B \rangle + \psi \langle A , C \rangle  
\end{equation*} 
The generalization of such concept is then defined to be the \textit{inner product}. We denote such with a few notations, namely $(\phi,\psi)$, the usual dot product will suffice, $\langle a,b \rangle $ or in a more quantum-ish notation, $\bra{a}\ket{b}$. The following are rules that it must obey.
\begin{align}
    \bra{a}\ket{b} &= \bra{b}\ket{a} \quad \text{(Symmetry)}\\ 
    \bra{a}\ket{a} & \geq 0 , \forall \ket{a} \neq \ket{0}\quad \text{(Positive semidefiniteness)}\\
    \bra{a} (\lambda\ket{W}+\psi\ket{c}) &\equiv \bra{a}\ket{\lambda b + \psi c} = \lambda\bra{a}\ket{b} + \psi\bra{a}\ket{c} \quad \text{(Linearity of the second argument)}
\end{align}
A vector space with an inner product is then called the \textbf{inner product space}\index{inner product space}. We note that the first argument of the inner product is often expressed to be \textit{antilinear}, that is $\bra{u}\ket{v}=(\bra{v}\ket{u})^{*}$ of the conjugate symmetry. In case of real number, since every real number equals its complex conjugate, then $\bra{u}\ket{v}=\bra{v}\ket{u}$ for $u,v\in E$. Any operation that satisfies this can be called the inner product. Hence, for example, we can define the inner product on the vector space of continuous real-valued function on the interval $[-1,1]$ by 
\begin{equation}
    \bra{f}\ket{g} = \int^{1}_{-1} f(x)g(x)\: dx
\end{equation}
of which then satisfies the inner product axioms. Continuing, we define a few concepts relevant to the space of inner product vector space. 
\begin{definition}[Orthorgonality]
    We say that two vectors $u,v\in E$ are orthogonal if $\bra{u}\ket{v}=0$. 
\end{definition}
\begin{definition}[Norm]
    The normal indicted on a vector space $E$ equipped with an inner product is defined as $||v||=\sqrt{\bra{v}\ket{v}}$ for $v\in E$. 
\end{definition}
\begin{definition}[Orthonormal basis]
    A set of basis vectors all of unit norm ($|u|=1$ for all $u$ of the basis), which are pairwise orthogonal, is called an \textit{orthonormal basis}. 
\end{definition}

The formula for the inner product, component-wise given a coordinate embedded space, is then as followed. Suppose we have two vectors, $\ket{V}$ and $\ket{W}$ expressed by the respective basis, 
\begin{equation}
    \ket{V} = \sum_{i} v_{i}\ket{i},\quad \ket{W} = \sum_{j} w_{j}\ket{j}
\end{equation}
we follow the axioms obeyed by the inner product to obtain 
\begin{equation}
    \bra{V}\ket{W}=\sum_{i}\sum_{j} v_{i}^{*}w_{j}\bra{i}\ket{j}
\end{equation}
where the $\bra{i}\ket{j}$ is the simple dot product between basis vector $i$ and $j$, and the term $v_{i}^{*}w_{j}$ is simply $v_{i}w_{j}$ if we work on real number. Why this equation? Simply because for Euclidean inner product on $\mathbb{F}^{n}$, the inner product formula is 
\begin{equation}
    \bra{(w_{1},w_{2},\dots,w_{n})} \ket{(v_{1},v_{2},\dots,v_{n})} = w_{1}v_{1}^{*} + \dots + w_{n}v^{*}_{n}
\end{equation}
However, there are structures of which this does not hold, hence the proportion is different. There, if $c_{1},\dots,c_{n}$ are positive numbers, of which act like weights, then an inner product can be defined as 
\begin{equation}
    \bra{(w_{1},w_{2},\dots,w_{n})} \ket{(v_{1},v_{2},\dots,v_{n})} = c_{1}w_{1}v_{1}^{*} + \dots + c_{n} w_{n}v^{*}_{n}
\end{equation}
\subsubsection{Gram-Schmidt scheme}

To go any further with the above inner product, we have to know the configuration $\bra{i}\ket{j}$, the inner product between basis vectors. That depends on the details of the basis vectors and all we know for sure is that they are linear independent (which is the criteria of a basis in an inner product space). Suppose that in a two-dimensional space where we have the basis vectors being linearly independent, but not orthogonal. Then, if we write all vectors in terms of this basis, compare to the unit default basis, then we have the inner dot product to be a double sum with four terms, two for each basis vector, as well as the vector components. However, if we use an orthonormal basis such as $\hat{i},\hat{j}$, only diagonal terms like $\bra{i}\ket{i}$ will survive, and we will get only the two terms, without overlapping. Hence that is why we want to have an orthogonal basis, or in general case, orthonormal basis. This is done using \textit{Gram-Schmidt scheme}.
\begin{theorem}[Gram-Schmidt]
    Given a linearly independent basis we can form the linear combinations of the basis vectors to obtain an orthonormal basis. 
\end{theorem}
The detail implementation of such is called the \textit{Gram-Schmidt procedure}, of which we take the virtue to not include here. The inner product then collapse to only $\bra{V}\ket{V}=\sum_{i}v^{*}_{i}w_{i}$, for such basis. 

\subsection{Dual spaces and Dirac notation}

Corresponding to any linear vector space $E$, there exists the dual space of linear functional on $E$. Such space is then called the vector space $E'$ of which for linear functional $F$ assigning a scalar $F(\phi)$ to each vector $phi$, such that to satisfy linearity criteria. To start, we need to know what is a linear functional. A \textit{linear functional}\index{linear functional} on a vector space $E$ is a linear mapping from $E$ to $\mathbb{F}$. From such, we can then define the collection of all such linear functional to a given structure the \textit{dual space}\index{dual space} of all linear functionals on $E$. In other words, $U'=\mathcal{L}(E,\mathbb{F})$. 

Normally, for a linear vector space, we can extrapolate some results and definitions.
\begin{theorem}[Dimensionality]
    Suppose $E$ is finite-dimensional. Then $E'$ is also finite dimensional and $\dim{(E')}=\dim{(E)}$.
\end{theorem} 

The dual space also has its own dual basis, and we will define it in similar fashion as followed (quite similar to a unit basis). 
\begin{definition}[Dual basis]
    If $v_{1},\dots,v_{n}$ is a basis of $E$, then the dual basis of $v_{1},\dots,v_{n}$ is the list $\varphi_{1},\dots,\varphi_{n}$ of elements of $E'$ such that each $\varphi_{j}$ is the linear functional on $E$ such that 
    \begin{equation}
        \varphi_{j}(v_{k}) = \begin{cases}
            1 & k = j \\
            0 & k \neq j
        \end{cases}
    \end{equation}
\end{definition}

The basis works the same for the dual space. Furthermore, we also have that the dual basis of a basis of $E$ consists of the linear functionals on $E$ that give the coefficient for expressing a vector in $E$ as a linear combination of the basis vectors. Quite a connection, but good nonetheless. We refer other concepts, like matrix dual, dual annihilator, dual space properties, dual maps, and so on, to our respective linear algebra chapters.

We note that it is indeed a scalar-inducing space via the inner product on the vector space $E$, if it is equipped with such operator. If $v\in E$, then the map that sends $u$ to $\langle u,v\rangle$ is considered as the linear functional on $E$. How do we know that this is true? Consider the inner product on the inner product space. Such operation is also one kind of linear functional, however, it is encoded into the vector space of such variation itself. What if we can find, for any linear functional in $E'$, an inner product representation of that specific expression? Fortunately, yes, and we leave the credit of such invention to Frigyes Riesz (1880-1956), who proved several theorems just like this, much earlier, and founded this particular representation theorem. Note that by default, linear functional is complete, such is to say they act on for every $u\in E$ of said vector space. 
\begin{theorem}[Riesz representation theorem]
    Suppose $V$ is finite-dimensional and $\varphi$ is a linear functional on $V$. Then, there exist a unique vector $v\in V$ such that 
    \begin{equation}
        \varphi(u) = \langle u,v\rangle
    \end{equation}
    for every $u\in V$. 
\end{theorem}
\begin{proof}
    We first need to show that there exists a vector $v\in V$ such that $\varphi(u)=\langle u,v\rangle$ for every $y\in V$. Let $e_{1},\dots,e_{n}$ be an orthonormal basis of $V$. Then, 
    \begin{align}
        \varphi(u) & = \varphi (\langle \langle u,e_{1} \rangle e_{1} + \dots + \langle  u,e_{n}\rangle e_{n} \rangle)\\
        & =  \langle u,e_{1} \rangle \varphi(e_{1}) + \dots + \langle  u,e_{n}\rangle \varphi(e_{n}) \\
        & = \left\langle u, \varphi^{*}(e_{1})e_{1} + \dots + \varphi^{*}(e_{n}) e_{n}\right\rangle
    \end{align}
    where $\varphi$ has the complex conjugate, also usually denoted as $\bar{\varphi}$. Nevertheless, if it is on real number, then the conjugate is just the real number. Thus setting 
    \begin{equation}
        v= \varphi^{*}(e_{1})e_{1} + \dots \varphi^{*}(e_{n})e_{n}
    \end{equation}
    we have $\varphi(u)=\langle u,v \rangle$ as desired. Now, we need to prove that only vector $v\in V$ has such behaviour. Suppose $v_{1},v_{2}\in V$ are such that 
    \begin{equation}
        \varphi(u) = \langle u,v_{1} \rangle = \langle u,v_{2} \rangle
    \end{equation}
    for every $u\in V$, then 
    \begin{equation}
        0 = \varphi(u) = \langle u,v_{1}\rangle - \langle u,v_{2} \rangle = \langle u,v_{1} - v_{2} \rangle
    \end{equation}
    for every $u\in V$. Taking $u=v_{1}-v_{2}$ shows that $v_{1}-v_{2}=0$. Thus $v_{1}=v_{2}$. This completes the proof of the uniqueness part of the result. 
\end{proof}
\subsection{Operators}
Operator is a step-up from your usual function. Let's say, if you are in Olympiad, probably heard of this in a different context being functional equations. Anyway, moving on, then, an \textit{operator} can be loosely defined as a rule like functions, but instead of variables, you give a function then produce another function outward. Such is to say, given an operator $A$, then $Af(x)=g(x)$ for $g(x)$ is called the \textit{product} of operator $A$ on $f(x)$. Notation varies, but usually it is the uppercase letter, or $\hat{A}$ for example. In physics specifically, of the 3D space, then we have $\hat{A}f(x,y,z,t)=g(x,y,z,t)$ as the functional product. 

As with a lot of things we are concerned of, we are interested in the class of \textit{linear operator} upon an (inner) vector space $V$. Sufficiently saying, denote the object of linear operator as $Q$, then if $\ket{\psi}$ is a vector living in such space, we have 
\begin{equation}
    Q(\alpha \ket{\psi}+ \beta \ket{\chi}) = \alpha(Q\ket{\psi}) + \beta (Q\ket{\chi})
\end{equation}
for $\alpha,\beta\in\mathbb{F}$. In context of quantum mechanics, we change from $\mathbb{F}$ to $\mathbb{C}$. 

To assert equality of two operators, $A=B$, we mean that $A\psi = B\psi$ for all vectors $\psi\in E$. Hence, there exists an isomorphism between them that is satisfied. We can then also define the basic composition rule for them, such as $(A+B)\psi = A\psi + B\psi$, and $AB\psi=A(B\psi)$, and so on. Because they are effectively a functional that acts on function, composition happens, and is written accordingly. For a scalar $k\in \mathbb{C}$, we also have $(kA)f=k(Af)$. Consequently, this means that 
\begin{equation}
    (\alpha \hat{A}+ \beta \hat{B})f(\cdot) = \alpha (\hat{A}f(\cdot)) + \beta 
\end{equation}
The operator can also be related toward the problem of linear algebraical solutions. For an operator $\hat{A}f=g$, there exists a special case that $\hat{A}f=\lambda f$, similar to the problem of eigenvalue and eigenvector. Thereby, instead of eigenvector, the main object of such space is the function space, we would have the concept of \textit{eigenfunction}\index{eigenfunction} $f$ and the eigenvalue $\lambda$ instead. In quantum mechanics, in general, we are interested in this case because with an operator of such kind, it is easier and interpretable of its components to convert from the state to the real space. How can we find the eigenvalue and eigenvector, though? Because we said that operators acts on the function space, then to examine such, we often have to observe how it changes a particular arbitrary function $f$. Let's take an example. 
\begin{example}
    To determine the eigenvalue and eigenfunction of $\hat{A}=i(d/dx)$, we have (suppose there exist such form of the eigenfunction):
    \begin{align}
        \hat{A}f &= \lambda f\\
        i \frac{d}{dx} f &= \lambda f\\
        \frac{d}{dx} f &= -i\lambda f
    \end{align}
    The above equation has the solution as $f=C\exp{(-i\lambda x)}$, and hence we see that $\lambda$ does not depend on $x$. Then there might exist infinitely many eigenvalues.
\end{example}

There are many examples of operators, in general, as their idea is sometimes just generalization of functions. Since they operate for mathematical operations on function, there exists some interesting examples to set. 

\begin{description}[leftmargin = 2.5cm]
    \item[Regular function] A regular function can be thought as  an operator, for $f: x\mapsto ax$, for $a\in \mathbb{R}$. The operator is $f$, operates on real numbers, and multiply the input $x$ by $a$. The function space then reduces to almost as \textit{scalar space}, and hence the function description fits. Squaring is also the same, as $f: x\mapsto x^{2}$. 
    \item[Functional application] For applying on functions, then is the proper terms for operator. $\hat{a}:f(x)\mapsto af(x)$ is an example, where operator $\hat{a}$ operates on scalar function that is mapped by multiplying by $a\in\mathbb{R}$.
    \item[Variable operator] The variable operator is the simple example, for $\hat{x}:f(x)\mapsto xf(x)$. This is the generalization of the above constant multiplication. 
    \item[Differential operator] The operator $D:f(x)\mapsto d/dx f(x)$ is the differential operator, operates on scalar function, which differentiates with respect to $x$ of such function. Partial differentiation operators, gradients, integral operators can be defined similarly. 
    \item[Momentum operator] The momentum operator $\hat{P}$ is defined by $\hat{P}= (\hbar/i)(\partial/\partial x)$, where $\hbar$ is the reduced Planck's constant, $i$ is the imaginary unit, of which operates on the wave function $\psi(x,t)$. For example, it can then be written as \begin{equation}
        \hat{P}(\psi) = \frac{\hbar}{i}\frac{\partial}{\partial x}:\psi(x,t)\mapsto\frac{\hbar}{i}\frac{\partial\psi(x,t)}{\partial x} 
    \end{equation}   
    \item[Energy operator] The energy operator is defined accordingly, as $\hat{E}=(\hbar/i)(\partial/\partial t)$, operates on the wave function. 
    \item[Hamiltonian operator] The Hamiltonian operator is defined, for $m$ is he mass of particle, regarded as constant, to be \begin{equation}
        \hat{H} = - \frac{\hbar^{2}}{2m} \frac{\partial^{2}}{\partial x^{2}}: \psi(x,t)\mapsto  \frac{\hbar^{2}}{2m} \frac{\partial^{2}\psi(x,t)}{\partial x^{2}}
    \end{equation} 
\end{description}

There is also the category of operators which act on vectors, though they are much more complex in notation (still matrix), but usually can be said of such for linear transformation in vector space. 

\subsection{Commutator}

Recalling the concept of an \textit{operator} $A: f^{(n)}(I)\to f(I)$ of which assigns to every function $f\in f^{(n)}(I)$ a function $A(f)\in f(I)$, of the mapping between two function spaces, then for any pair of operators, $A, B$, the \textbf{commutator}\index{commutator} of them is defined as 
\begin{equation}
    \left[ A,B \right] = AB - BA
\end{equation}
The notation here varies, but usually, it is used with the bracket $\left[\cdot\right]$, for operator notation either $A$ or $\hat{A}$ for example (to not mix up stuffs). Simple properties of commutator on operators can be as followed. 

\begin{theorem}[Commutator properties]
    Let $a,b$ be constants, $A,B$ be operators on $n$-differentiable space. Then the following set of equalities hold: 
    \begin{align}
        [f(x),x] & = 0\\
        [A,A] & = 0 \\
        [A,B] & = - [B,A]\\
        [A,BC] & = [A,C]B + A[B,C]\\
        [AB,C] & = [A,C]B + A [B, C]\\
        [a+ A, b+ B] & = [A, B]\\
        [A+B, C+D] &= [A,C] + [A,D] + [B,C] + [B,D]
    \end{align}
\end{theorem}
In general, you can think of the commutator as the interference measurement device. As least to the theory of operator, it transforms the function space, or in general can also be the scalar space, with different ordering of the transformation applied. Remember that operator is indeed a mapping. For case that operator transform on scalar space, for example, $\mathbb{F}^{n}$, then operator reduces to function, and for function space, it maps from this function space to another, basically in both case it works as a transformation on the base space. In general, it basically measures the difference between $AB$, apply $A$ then $B$ of which both acts and transform the space $X\to Y$, then $Y\to Z$, or $BA$, which is the reverse so $X\to Z$ then $Z\to Y$. If they are the same, so basically circular to a result, then the commutator measures zero. This is seen as $[A,A]=0$, as you apply $A$, then continue applying $A$, or by the reverse the path is the same (as connecting two straight lines). We can also then specifically see that it is very much similar to the theory of group, and indeed can be generalized in such direction. The commutator of two group element $A$ and $B$ is $AB^{-1}A^{-1}B^{-1}$, and two elements $A$ and $B$ are said to commute when their commutator is the identity element (same as us). When the group is the \textit{Lie group}, the Lie bracket in its Lie algebra is an infinitesimal version of the group commutator. We defer this to a separate issue on its own, but the idea is the same. 

The last identities to think of aside from those includes the Jacobi's identity, 
\begin{equation}
    [A,[B,C]] + [B,[C,A]] + [C,[A,B]] = 0
\end{equation}
and the generalization of nested operators $[A^n, B]$ etc, 
\begin{equation}
    \left[A,B^{n}\right] = \sum^{n-1}_{j=0}B^{j}[A,B]B^{n-j-1}, \quad \left[A^{n},B\right] = \sum^{n-1}_{j=0} A^{n-j-1}[A,B]A^{j}
\end{equation}
In quantum physics, in general, such expression hence evaluates the interference between two properties of the physical system. If they simply commute, then it is guaranteed for both of them to be measurable on the state space of $\ket{\psi}$ as the wavefunction. 

\subsection{Hermitian operators}

To understand Hermitian operators, we need to understand the adjoints of operators. For now, we are solely in the context of the complex field $\mathbb{C}$, so any kind of adjoint would be between complex operator. By convention, any operator on such function vector space, such as linear operators, can be represented by matrices. Let $H_{1}$ and $H_{1}$ be Hilbert space $X$ over a field $\mathbb{F}\in\{\mathbb{R},\mathbb{C}\}$. We define the adjoint\index{adjoint operator} as followed:
\begin{theorem}[Adjoint operator]
    Let $T\in\mathcal{B}(X)$ be a bounded linear operator on a Hilbert space $X$. There exists a unique operator $T^{*}\in\mathcal{B}(X)$ such that 
    \begin{equation}
        \langle T x, y\rangle = \langle x, T^{*}y\rangle, \quad \forall x,y\in X
    \end{equation}
    The operator $T^{*}$ is called the \textbf{adjoint} of $T$. 
\end{theorem}
This is true for also finite dimensional inner product space, and so is Hilbert space through assumptions on Riesz representation theorem. The notion of adjoint is the functional generalization to the concept of transpose. By the definition of transpose, a transpose on a matrix $X$ as a linear transformation reverses the direction of said transformation on the inner product space per definition of said space. In a sense, you can think of it as transforming perspective, from linear operations on $\mathbb{V}$ into the space of covectors, or dual space $\mathbb{V}'$. Remember that the dual space represents measurements and functional application on vector space. Hence, basically, transpose do the operations on $W^{*}\to V^{*}$ for a linear transformation $V\to W$. This works on the complex vector space, where the inner product involves complex conjugation, and in such case there exists the \textbf{adjoint matrix} for a given complex transformation. 

Something is then called a \textbf{Hermitian operator} if the adjoint is the operator, or $T^{*}=T$. Quantum mechanically, for $\psi, \phi$ being arbitrary functions on complex space. Then, if operator $\hat{A}$ is such that \index{Hermitian operator}
\begin{equation}
    \int \psi (\hat{A}\phi)^{*}\: d\tau = \in t \psi^{*}(\hat{A}\psi)\: d\tau
\end{equation}
then it is called a quantum Hermitian operator, or \textit{self-adjoint} operator. Under matrix notation, then adjoint is the \textit{conjugate transpose} of such complex matrix. We have the following definition for \textit{skew-adjoint} and \textit{self-adjoint} operators:
\begin{definition}[Skew-adjoint operators]\index{skew-adjoint}
    A linear operator $A: \mathcal{U}\to\mathcal{Q}$ is said to be \textit{skew-adjoint} if $A^{*}=-A$, that is 
    \begin{equation}
        \langle A\mathbf{u}, \mathbf{v}\rangle = -\langle \mathbf{u} , A\mathbf{V}\rangle
    \end{equation}
    for all $\mathbf{u,v}\in \mathcal{U}$. 
\end{definition}
\begin{definition}[Self-adjoint operators]\index{self-adjoint}
    A linear operator $A: \mathcal{U}\to\mathcal{Q}$ is said to be \textit{self-adjoint} if $A^{*}=A$, that is 
    \begin{equation}
        \langle A\mathbf{u}, \mathbf{v}\rangle = \langle \mathbf{u} , A\mathbf{V}\rangle
    \end{equation}
    for all $\mathbf{u,v}\in \mathcal{U}$. 
\end{definition}
All of these definitions are applicable for both real and complex space, of the metric standard. 
\subsubsection{A closer look at complex adjoint}

Previously, we kind of mix up between the field $\mathbb{F}$, where it is arbitrary of the complex or real space. For now, let us examine such adjoint in the complex space itself. 

Given the linear mapping $A:\mathbb{C}^{n}\to \mathbb{C}^{m}$, then it is easiest to think of 
\begin{equation}
    A\mathbf{u} = u_{1} \mathbf{a}_{1} + \dots u_{n}\mathbf{a}_{n}
\end{equation}
for the transformation basis $\mathbf{a}_{i}$. Thus, we may deduce the $i$th entry of $A\mathbf{u}$ given by 
\begin{equation}
    (A\mathbf{u})_{i} = u_{1}(\mathbf{a}_{1})_{i}+ \dots u_{n}(\mathbf{a}_{n})_{i} = \sum^{n}_{j=1}u_{j}a_{ij}
\end{equation}
We define the complex version of the inner product as the \textbf{Hermitian inner product}, such as followed. 
\begin{definition}[Complex inner product]
    The vector space $V$ over field $\mathbb{C}$ is said to have the complex inner product $\langle\cdot,\cdot\rangle:V\times V\to \mathbb{C}$ if and only if it satisfies, for all $x,y,y_{1},y_{2}\in V$ and $\alpha\in \mathbb{C}$ the following set of conditions:
    \begin{align}
        \langle x, y_{1}+y_{2}\rangle &= \langle x,y_{1}\rangle + \langle x, y_{2}\rangle \\
        \langle x, \alpha y\rangle & = \alpha \langle x,y \rangle\\
        \langle y, x \rangle &= \langle x,y\rangle^{*}\\
        \langle x, x \rangle &> 0 \: \text{if}\: x\ne 0
    \end{align}
\end{definition}

Thus, if $(A\mathbf{u})_{i} =\sum^{n}_{j=1}u_{j}a_{ij}$, then we have the inner product as 
\begin{equation}
    \begin{split}
        \langle A\mathbf{u},\mathbf{v} \rangle 
        & = \sum^{m}_{i=1} \left(\sum^{n}_{j=1}(u_{j}a_{ij})\right)^{*}v_{i}\\
        & = \sum^{m}_{i=1} \left(\sum^{n}_{j=1}(u_{j}^{*}a_{ij}^{*})v_{i}\right)\\
        & = \sum^{n}_{j=1}\left(\sum^{m}_{i=1}(u_{j}^{*}a_{ij}^{*})v_{i}\right)\\
        & = \sum^{n}_{j=1}\left(\sum^{m}_{i=1}(v_{i}a_{ij}^{*})u^{*}_{j}\right)\\
        & = \sum^{n}_{j=1}\left(\sum^{m}_{i=1}(v_{i}a_{ij}^{*})\right)u^{*}_{j} \\
        & = \langle \mathbf{u}, A^{*}\mathbf{v}\rangle
    \end{split}
\end{equation}

Here, the $ji$ entry of $A^{*}$ is the conjugate of the $ij$ entry of $A$. How do we interpret such result? Basically, here we gain \begin{equation}
    (A^{*}\mathbf{v})_{i} = \sum^{m}_{i=1} (v_{i}a_{ij}^{*}), \quad j = 1,\dots, n
\end{equation}
Of which looks like the dual form of $(A\mathbf{u})_{i}$. And indeed it is so much of the change in basis, especially when we have the decomposition: 
\begin{equation}
    (A^{*}\mathbf{v}) = \sum^{m}_{i} v_{i}a^{*}_{ij} = v_{1} a^{*}_{1j} + \dots v_{m}a_{mj}^{*} 
\end{equation}
This is the linear combination of the conjugates of the rows of $A$. In general, if a theorem refers to the adjoint of a linear map, then if you are applying that to linear maps on finite dimensional complex vector spaces, you can equivalently substitute that with the phrase conjugate transpose. The conjugate transpose interact with the inner product as the general adjoint is. For a linear map $A:\mathbb{C}^{n}\to\mathbb{C}^{m}$, then $A$ can be represented as an $m\times n$ complex matrix. The adjoint by then is the $n\times m$ matrix that is the matrix $A$ reflected through the diagonal with each entry conjugated, and we call it conjugate transpose. 
\subsubsection{Quantum Hermitian operators}
As we have said, the operator $P$ is defined as Hermitian if it satisfies 
\begin{equation}
    \int \psi^{*}(x) P \phi(x) \: dx = \int (P \psi(x))^{*}\phi(x)\: dx
\end{equation}
To see how this definition works, let's show that the momentum operator, $p=-\hbar d/dx$ is Hermitian. We have: 
\begin{align}
    \int \psi^{*}(x) \left(-i\hbar \frac{d}{dx}\right)\phi(x)\: dx 
    & = -i\hbar \int \psi^{*}(x)\frac{d}{dx}\phi(x)\: dx \\
    & = -i\hbar \int \psi(x)^{*} d\phi(x)\\
    & = -i\hbar \left\{ \psi^{*}(x)\phi(x) - \int \phi(x)d\psi^{*}(x) \right\}\\
    & = \int \phi(x)\left(-i\hbar \frac{d}{dx}\psi(x)\right)^{*}
\end{align}
which is a Hermitian by definition of hermiticity. There are three important consequences of an operator being hermitian. Its eigenvalues are \textit{real}, its eigenfunction corresponding to different eigenvalues are orthogonal to each other, and the set of all its eigenfunctions is complete. Observations in reality is always measured of real number, hence we can interpret the eigenvalues as observable quantities of a given quantum system (we will get to this seriously later on). Eigenfunctions being orthogonal basically allows us, to get exclusive state. Hence, across all states, you can write a general state $\Phi$ as the combination of the orthogonal states as 
\begin{equation}
    \Phi = c_{1}\phi_{1} + c_{2}\phi_{2} + \dots 
\end{equation}
for the $\{\phi_{i}\}$ combinations. Completeness basically means that eigenfunctions here form a complete basis of the state space, such that any $\Phi$ can be represented with just as many combinations. 

\clearpage 

\section{Theoretical alternatives - synthetic quantum-like}


\clearpage

\section{The quantum well structures case}

Quantum well, or rather, specified topology on coordinate space of specific mechanistic landscapes, has been one of the central part in the framework of quantum mechanics. Indeed, foundational sources, textbooks and materials in quantum physics considers different forms of such quantum confinement environment as to gauge how a non-free space system evolutes and changes, which align with much of the observable results. While textbooks and works \cite{griffiths2018,sakurai2017,landau1981} has been introducing the concepts and working of quantum potential confinement and so on, the inner working of systems in practical setting or of more complicated requirements, such as semiconductors and typical design space, as seen of historical works or textbooks \cite{kane1957,chang1980,kittel2005,bastard1988,davies1998,harrison2016,weisbuch1991,haugkoch2009,yu2010,barsan2016,esaki1974}, and require much more sophisticated mechanisms or theoretical tooling to work. As such, we are inclined to provide a partial analysis of target, on the structure and treatment of \textit{quantum confinement structures}, and the deliberate exact analytical results and limitation in forcing such approach, if we are to be technical and of correct treatments to the quantum discourse. 

Quantum well structure branched out from quantum theory, and the typical potential insertion as a more practical, or more physically relevant model of considering real system beyond finite and infinite well. As such, currently, they are more correlated to be defined as \textit{thin layered semiconductor structures} in which we can observe and control quantum mechanical effects, all of which are non-free particle system of which can be naturally said of to be the ambiance \textbf{topology} of the physical system itself (per Schrödinger equation expression and so on). As such, per David A. B. Miller \cite{David2020OpticalPO}, the quantum well structure and its relevant class of restriction (quantum wire and quantum dot) derive most of their special properties from "the quantum confinement of charge carriers in thin layers, of one semiconductor well material sandwiched between other semiconductor barrier layers". If the middle layer is effectively non-impeding in any sense, be it vacuum or transmittable medium, the topology is analogous to square potential restriction, of which is analogous to quantum potential models in theory. They are useful, as laboratory setting in which can be used to set up much more complicated structures or setting, with feasibility of mechanism control than directly observe, for example, excitons at large. In essence, we can say that quantum well or potential-induced structures, makes up the bulk of relevant topology-creating physical system quantum theory most often consider itself in of scale; and simply speak, quantum theory becomes physically meaningful only once a potential defines structural topology to work on it. The easiest justification of such, is from the observation of the fact that in the current model, all nontrivial quantum topology enters through constraints on the Hamiltonian.

We would like to clarify the usage of \textbf{topology} here is rather non-standard, and would clash with typical intuition on algebraic or differential topology, and so on. The usage here is rather \textbf{ontological} in its core argument, and it puts the physical reality on the \textbf{first-order} basis (highest). The need for this clarification is justified, as for even sophisticated reader, and a lot of Large Language Model (LLM) of existence, failed to read this part, or discard it as unimportant directly, and criticize it later on. In a sense, automated machines or formalist readers systematically discard ontological framing and then retroactively impose mathematical standards, which is what we are to prevent our interpretation to be of. Thus, while sharing the naming, rather, what we want to emphasize here is the definition on which hinges on the onset of, for example in this relevant context, \textit{how the Hamiltonian carves out configuration space and Hilbert space} into allowed and forbidden regions, or rather, what exists before a theory, or treatment, can talk about representation on smoothness, manifolds, or bundles, that are justified or not. In such sense, topology here is \textbf{first-order}, basing the system in consideration of all admissible solutions first and foremost, and is created from confinement, boundary, or periodic conditions, and domain restriction of operators. It is then closer to spectral topology, or operator-domain topology than to differential topology, of which would be considered in such case to be the \textbf{second-order topology}, induced by mathematical structures in interpretation of the configuration space. Naturally, the usage would be fairly standard once it is clarified. In Schrödinger theory or to the larger degree of quantum theory, there are normally three structural inputs that shape the system. First is the base space, which is usually $\mathbb{R}^{n}$, the Hamiltonian operator, and the domain on which that operator is self-adjoint to be of interest. As such, in the expression of the Schrödinger equation, $V(x)$ is the most important external fixture of which applies on all three. Dirac equation inherits the same topology, but with added first or second-order topology only, for example, the \textbf{spinor topology}, relativistic dispersion, or gauge coupling as geometric potential. In practice, especially for structure of interest, Dirac equation is reduced back to Schrodinger-type envelope equations via nonrelativistic limits or Foldy-Wouthuysen transformations, which reduces the equation to Schrödinger-Pauli equation. In fact, by observation, the Schrödinger equation,
\begin{equation}
    i\hbar \partial_{t}\psi = \left(-\frac{\hbar^2}{2m}\nabla^{2}+V(x)\right)\psi
\end{equation}
and the Dirac equation on relativistic formulation,
\begin{equation}
    i\hbar \partial_t \Psi = (-i\hbar \bm{\alpha}\cdot \nabla + \beta mc^{2}+ V(x))\Psi
\end{equation}
share the same $V(x)$ external topology generation. Naturally so, we simply mean topology here as \textit{the constraint-induced structure of the quantum state space}. 

While the infinite well is a mathematical idealisation (hard walls, zero penetration); and the idealized theoretical finite well adds tunnelling and leakage, which is physically relevant when barriers are comparable to particle energy; the practical picture is rather much more different and complex, specifically considered of the potential resides inside the well, $V(x)$, and the shape of the barrier with admissible interaction. As such, the creation of models and explicit structures that can be replicated of such potential device, hinges on the goal for examining those conditions of the quantum system where their shapes are not trivial, and can induce complex interaction patterns, for example, \textit{nonlinear absorption} inside crystal lattices, all of which can be practically observed, or made to be near-approximation of ideal cases. From such goal, we also discovered, or rather made use of many contingent structures or interesting ordering beyond simply a single-potential system (in reality, such confinement in singularity is rather hard to make), for example, replacing isolated wells by a periodic atomic array, which leads to the concept of \textit{band formation} in such specific structural containment. The earliest example of this is from Kronig-Penney model \cite{10.1098/rspa.1931.0019}, of which studies the square-well periodic potential system. Such system considers
\begin{equation}
    V(x) = \begin{cases}
        V_{0} & -b \leq x \leq 0 \\
        0 & 0 < x \leq a
    \end{cases}
    , \quad 
    \psi(x+c) = \beta\psi(x)
\end{equation}

where $\beta$ is the constant of magnitude unity, or $\beta=\exp{(2\pi i l /n)}$, for $l=0,1,\dots,n-1$, and $c=a+b$ as the periodic length of the potential. The periodic property of the model can also be expressed as
\begin{equation}
    \psi(x+c) = e^{i\kappa c}\psi(x), \quad \kappa = 2\pi l/nc
\end{equation}
with nondimensionalization (for more, see \cite{li2024onedimensionperiodicpotentialsschrodinger} for also a fairly modern way of solving such system via finite difference method). Overall, the formulation started the conceptual bridge from single well to bulk semiconductor electronic structure, and thus enable the further road to actualize such well, and extending the topological formation system into much more active and sophisticated platforms. 

\subsection{Conductive and semiconductor}

The realization of quantum well, and later low-dimensional devices being created from conductive-capable or semiconductive materials, is rather an important step. Shifts to designing more sophisticated structures and topology is influenced by the capacity to actualize them on true physical devices, and hence, we are to give requirement to examine how the devices can form capture topology, and how they are used. Naturally, it also is not a one-way road, as understanding confining potential well topology helps in analysing well structures and material properties, and thus design philosophy. For a simple example, such is why there has been surges in the early 2000s for theoretical analysis on low-dimensional confining analysis, to analyse different properties of \textbf{Graphite}-based devices, mostly within nonlinear absorption phenomena, where conductivity of semiconductors can vary with environmental parameters, like temperature, light or electric current, with specifically important application as \textbf{optical computing}, and \textbf{optoelectronic devices}. 

It is then our task to provide a quick overview on how semiconductor device works, and the process in obtaining our quantum well confinement devices. Within such, we will develop an intuition on how the system is modelled, and to befit our experiences afterward. 

Semiconductor is constructed from solid-state material and structures. As said the topology of such material, for example, crystalline or at least solid lattice, creates \textbf{electronic behaviours} from electrons occupying energy bands inside, which are much closer than normal discretized energy level \cite{huang2DSemiconductorsSpecific2022}. Thence, at least for now, properties of solid, including \textit{doping} on shifted Fermi level, mechanical, thermal, fabrication scalability and stability, are what driving the usage of semiconductor in solid. While organic and molecular semiconductors \cite{FACCHETTI200728}, molecular electronics, and iontronics or bio-ionic devices are in rather active research, they cannot yet replace semiconductor for solid-state physics. Hence, for now, we would be fixed to this. 

Solid-states materials can be categorized into three classes - insulators, semiconductors, and conductors. They are categorized based on resistivity and conductivity, thus there are usually no fixed choice, for example as hydrogen being the standard baseline of molecular structure (as most analysis focused on them), but rather on what kind of compounds or element fits onto such category. Typical fit for semiconductor is with binary III-V columns semiconductor, or silicon and germanium single-element semiconductor. For semiconductor specifically, separation of bands happens when there exists more than two atoms together in interaction. Take Bohr model for example. The energy levels of a hydrogen atom are given by
\begin{equation}
    E_{H} = \frac{m_{o}q^{4}}{8\epsilon_{o}^{2}h^{2}n^2}
\end{equation}
for $m_{o}$ the free electron mass, $q$ the electron charge, $\epsilon_{o}$ the free space permittivity, $n$ the principle quantum number of the atom, and $h$ the Planck constant. By Heitler-London picture, when two hydrogens atoms $A$ and $B$ approach, each isolated \texttt{1s} level (note that this is much more general as for \textit{all electron in interaction}, hydrogen is the special case), $\ket{\phi_{A}}$ and $\ket{\phi_{B}}$ interacts through the two-electron Hamiltonian,
\begin{equation}
    \hat{H} = \hat{H}_{A} + \hat{H}_{B} + \hat{V}_{AB}
\end{equation}
for $\hat{V}_{AB}$ the Coulomb interactions between electrons and between electrons and the far nucleus. The molecular states specified thence reduces to
\begin{equation}
    \ket{\phi_{\pm}} = \frac{1}{\sqrt{2(1\pm \braket{\phi_{A}|\phi_{B}})}}(\ket{\phi_{A}\phi_{B}}\pm\phi_{B}\phi_{A})
\end{equation}
and solving this gives us the energy level description,
\begin{equation}
    E_{\pm} = \frac{E_{0}\pm H_{AB}}{1\pm \braket{\phi_{A}|\phi_{B}}}
\end{equation}
where $E_{0}$ is now the isolated atomic level, and $H_{AB}$ is expressed into the interaction matrix element. This allows the splitting of energy level in hydrogen, which is
\begin{equation}
    E_{+} - E_{-} \approx 2H_{AB} = 2 \braket{\phi_{A}|\hat{H}|\phi_{B}}
\end{equation}
This result is the basis for both Heitler-London and Slater-Pauling molecular orbital theory, and thus precedent to explain the splitting of such. For historical primary sources, see Heitler-London paper \cite{HeitlerLondon1927}, Slater and Koster's paper \cite{SlaterKoster1954} for Slater-Koster parameterization, and development, derivation of such results \cite{Papaconstantopoulos2003,Harrison1980,AshcroftMermin1976,kittel2005,SzaboOstlund1989,KolosWolniewicz1965}. 

As seen from such, $E_{H}$ of the energy band level for hydrogen sets the reference level $E_{0}$, i.e. the absolute bound-state energies of an isolated hydrogen atom \footnote{Note that the reference level is a bookkeeping device, not necessarily a choice for minimum energy, while there are cases where they coincide. Another confusion point is that typically, the scale is inverted. That is, we set $E(\infty)=0$, thus for all minimal ground state of actionable and physical system, the baseline is mathematically encoded as negative. The choice of Hydrogen is because it gives ideally analytical non-relativistic Coulomb bound system, which is better than having a chaotic system with non-trivial properties as reference point. Since energy of a system is encoded by the Hamiltonian, typically the definition of ground state, would be $\min{\sigma(H)}$ for the minima of all possible energy values of the Hamiltonian, assuming discrete quantization Hamiltonian operator.}, and splitting can be explained in a populated system due to interaction enters through the off-diagonal matrix element $H_{AB}$. Within this mathematical representation for the splitting, such is what drives the separation. Extending such to crystal requires only treating each bound atomic level $E_{0}$ of an atom becoming a set of many interacting sites. This is called \textbf{tight banding}, and said can be expressed via the tight-binding structure, 
\begin{equation}
    E(\mathbf{k}) = E_{0} + \sum_{\mathbf{R}\neq 0} t(\mathbf{R})e^{ik\cdot \mathbf{R}}
\end{equation}
for $E(\mathbf{k})$ the energy eigenvalue for crystal momentum $\mathbf{k}$-space, constrained by crystal translational symmetry, $\mathbf{R}$ is the lattice vector for atomic sites, $E_{0}$ the on-site energy (atom without hopping anywhere), $t(\mathbf{R})$ is the hopping integral, $e^{i\mathbf{k}\cdot \mathbf{R}}$ the phase factor from translational symmetry topology, and $\phi_{0},\phi_{\mathbf{R}}$ are the two localized orbital by origin and translation by $\mathbf{R}$ of the LCAO basis. One thing to note about such scheme is that, translational symmetry is directly the potential encoded structure itself. Such is why it constraints the available configuration for $\mathbf{k}$ crystal eigenstates, as distinct representation of the translation symmetry condition, with stability, and only those symmetry-compatible states can exist as stable stationary states. Such is also partly why there exists the \textbf{Brillouin zone}, which is the \textbf{minimal representation} of the reciprocal lattice vector group. In fact, you do not need Brillouin at all, and devise the extended reciprocal space with only sufficient counting and identification scheme, and nothing would change physically. Such is also the idea of a physical theory being gauge, and such direction of thought is popular in modern physics. 

Now, since we have energy bands structures, such system can be further elaborated to finally explicitly gives a more macroscopic version relative to such internal crystalline structure itself, and producing non-trivial, stable quantum well, which fits both being used as a fabricated experiment, and many electronic behaviours. Specifically, in theory, we posit the potential shape $V(x)$ and solve the Schrödinger equation for energies in such space. Here, we attempt to stack different material with different band edges together, producing the well itself with similar topology to the proposed theoretical solution. Because from the dispersion $E(\mathbf{k})$, for each band, there is a maximum and minimum, which would be called \textbf{edges} of the band. Two types of bands are available for semiconductor device, of which are not in metal, for example. 
\begin{enumerate}[itemsep=1pt,topsep=1pt]
    \item \textbf{Conductive band}: where electrons are mobile because of the existence of the lowest-energy band that is empty, above the Fermi level. 
    \item \textbf{Valance band}: where electrons are bound in bonds, of the highest-energy band that is totally filled. 
\end{enumerate}
The split comes from electron occupancy, which allows you to do doping, temperature morph, and external fields enforcement. For now, however, we return to the stacking material. Because of the existence of bands, stacking different bands together creates \textbf{broken periodicity} inside the crystalline complex structure. Periodicity forbids confinement from translational symmetry, hence effectively, this makes a continuous quantum well, where conduction or valence band edge varies continuously in space. The spatial variation is the potential well $V(z)$ that we want, and it contains the band profile in full. 

At this point, there exist two topologies. First, is the physical material device topology of the stacked lattice structure, basically the setting constrained physical space. Second, is the topology of the reduced Hilbert space on what we model. Conductive band and valance band gives us two scenarios. Undoped semiconductor at zero temperature follows the description of valance band having \textit{full electron}, conduction band with \textit{empty electron sea}, and there exists no holes in the valance band. When we $p$-doping the occupancy, removing some electrons from the valence band, we create missing holes, similar to how optical excitation of the same result. On the conductive band, we \textit{add electron} to the topological configuration itself. The two scenarios can now be realized. First, is the single-band electron topology, where we project onto one conduction band, and use EMA to capture its behaviours. Second is the single-band hole topology, where we project onto valence band, and capture hole quantum wells, heavy-hole or light-hole scenario, anisotropic effective mass, and $p$-type devices. Such can also be extended, if one wishes, for lower-dimensional devices. What we have done, however, can be said to be \textbf{physical simulation} of type, where the device itself simulates the Hamiltonian. 

Certain analysis can be made at this stage. Mainly, is the requirement for the expected value picture here. Implicitly, Bloch's theorem and lattice structure formation works within the boundary of \textit{multi-particle system}, and thus cannot be introduced into a single isolated condition setting. We can, indeed, operationally ignore the full valance/conduction band structure after we fixed their degeneracy and symmetry of the testbed CBM/VBM, discarding irrelevant high-energy band, and keep the slow envelop dynamics or so, thus allowing for the creation of 1D bound-state problem. Furthermore, we can use such testbed and doping mechanism to fabricate single-point quantum dots, or single-hole traps, charge qubits, and so on. However, such picture assumes we already have the multi-particle, ensemble system on itself, since Bloch bands and valence condition descriptions only make sense as many-particle, many-electron constructions. Furthermore, practical physics for said devices require us to give expectation i.e. many-body approximation, because of uncertainty and often the perspective of \textit{Fermi sea transition} instead of singular unstable determination. As such, before reducing the picture into single particle and fabrication of such singularity, we are required to have the devices. Which themselves are not single. The singular picture still works, however, because of the external topology existence. From the point of view of the single-particle, we effectively treat a large number of degrees of freedom to be absorbed into the external structures that the single particle moves in. That is why, even with single electron potential creation, and the process to confine such, the single-electron view still works, though with support from many other components. For more technical and correct creation of well system, we can defer to authoritative source on such matter \cite{David2020OpticalPO,davies1998}. 

\subsubsection{Low-dimensional semiconductor}

From such formation, low-dimensional consideration is a natural branching into orthogonal aspect of practical systems. The term low-dimensional system, and its realization to semiconductor, stems from the observation that for a full dimensional setting, adding structures equal to reducing free-space transport into lower dimensional freedom, i.e. applying potential topology effective on different spatial domain, thus reducing the subject into quantum well (one-dimensional restriction), quantum rod (two-dimensional restriction) or quantum dot (three-dimensional restriction). Peeters \& Oscar \cite{Peeters1992LowDS} introduced such feasibility from the application of molecular beam epitaxy (MBE), to create heterojunction structures and quantum devices. Again, we defer to authoritative overview on such matters, however, from Peeters himself, such is said about the creation of low-dimensional devices. 

\begin{figure}[htb]
    \includegraphics[width=\linewidth]{img/img1.png}
    \caption{Schematic of low-dimensional semiconductor structures and the schema of energy level, with respect to the density of states $D(E)$.}
\end{figure}

\begin{quote}
    Continuing advances in MBE and breakthroughs in nanolithography has made it possible to confine the electrons; further in more dimensions which leads to quasi-one and quasi-zero dimensional semiconductor structure. Quasi-zero dimensional structures behave like artificial atoms. For example quasi-one-dimensional MOSFET in silicon have revealed the universal conductance fluctuations due to random quantum interference predicted by Altshuler et al. \cite{Peeters1992LowDS}.
\end{quote}
He also emphasized that the essentiality of low-dimensional systems is their particular density of states. The latter in which was mentioned, is "the number of available electron states per unit energy", all of which links to the properties of transport and confinement protocols.

\subsection{Structural commitments}
When moving from textbook wells to realistic condensed-matter heterostructures, several modelling approximations and their domains of validity become critical, especially considering physical realization of such. First, is the assumption on non-interacting electron inside confinement. Such is to say, we neglect electron-electron correlation and excitonic binding unless explicitly included in the model. This basically gives us formulaic strength - to solve a many-body system of well-behaved independency, we can solve only the system of Schrödinger equation with single-particle formulation, thus staying in the zone of the first quantization effectively. The method works practically well as an \textit{approximate simplification} condition for various cases, and is a requirement for the working of the free electron model, nearly-free electron model, alongside Bloch's theorem for crystals lattice, as seen in literature \cite{Fulde1995,McGuire_1997}. However, in region where this assumption might be the dampening prediction, or where such electron-electron interaction cannot be considered null as for them to be the major mechanism, divergence happens, often strong breakdown, as seen in particular, of the sequential double ionization problem space \cite{Pfeiffer_2011}. 

Second, is the effective-mass approximation (EMA) usage of the carrier electron, thus also $k\cdot p$ expansion. Per historical development, we get a model of Bloch's electron (1929) from Bloch's theorem, in which dynamics in a periodic potential can be expressed in such. In principle, Bloch's wave theorem for specifying electron form, similar to the independent electron approximation, assumes that such many-electron interaction is lumped together into an effective electron potential. Doing such, it dramatically reduces the complexity of the model, and thus having him to solve only the reduced setting for
\begin{equation}
    \begin{split}
        &\hat{H}\psi(\mathbf{r}) = \left(-\frac{\hbar^2}{2m}\nabla^{2} + V(\mathbf{r})\right)\psi(\mathbf{r}) = E\psi(\mathbf{r})\\
        & \Delta \psi + \mu (E-V)\psi = 0, \quad \mu = (8\pi^{2}m)/\hbar^{2}
    \end{split}
\end{equation}
for $V(\mathbf{r})$ is a periodic potential with the lattice periodicity as said before. On the space of three-dimensional system consideration of the crystal, the solution that Bloch reached from his paper \cite{blochUberQuantenmechanikElektronen1929a}, formulated the solution of the wavefunction as
\begin{equation}
    \psi_{klm}(x,y,z) = \exp{\left( 2\pi i \dfrac{kx}{K}+\dfrac{ky}{L}+\dfrac{kz}{M} \right)}u_{klm}(x,y,z)
\end{equation}
More details can be taken from the original paper, or review papers of the original. Per simplification models, as specified of Rabinovitch \cite{PhysRevB.4.1017}, Hyeonwoo et al. \cite{Yeo_2020}, or Donetti et al. \cite{10.1063/1.3639281}, EMA is a very effective, simplified model of Bloch's electron under certain conditions, for both computational and numerical simulation purpose, as well as analytical one for well-behaved system. For a nearly parabolic band near extremum $k_{0}$, the effective mass tensor $m_{ij}^{*}$, energy dispersion $E(\mathbf{k})$, $\mathbf{k}$ being the crystal wavevector can be connected as
\begin{equation}
    \frac{1}{m_{ij}^{*}} = \frac{1}{\hbar^2} \frac{\partial^{2}}{\partial k_{i}\partial k_{j}}E(\mathbf{k}) \Bigg\rvert_{\mathbf{k}=k_{0}}
\end{equation}
However, as they also pointed out, the method can be considered crude approximation in many settings, for example, hole transport, where "strong non-parabolicity and anisotropy of silicon valance band makes EMA a very crude approximation". Alongside envelop function \footnote{The envelope function refers to a mathematical representation of a wavefunction in quantum wells, characterized by a plane wave in the in-plane directions (x, y) and a localized wavefunction in the z-direction, normalized to the sample area. It describes the spatial distribution of particles confined in a quantum system while allowing for extended states in the unconfined directions \cite{EnvelopeFunctionOverview}.} and $k\cdot p$-expansion \cite{Li_2014} (\cite{M_G_Burt_1999} for fundamentals discussion of both conventional and envelop theories for modern semiconductor physics of the time), they form the bulk of theoretical tools for working on slow-varying spatial modulation of the electron, in other word, when the expressed envelop representation varies slowly on the scale of a lattice constant and inter-band mixing is small. As such, it breaks down for also regions with abrupt interfaces, strong strains, or very narrow wells. The representation of envelope function for heterojunctions also introduces boundary conditions that are often unsafe and should be proceeded carefully, and thus requires in certain case, justifications, as seen for microstructures \cite{Burt1992TheJF}. 

Along the way, there are practical assumptions on smoothness, and penetration depth to be enough for quantization. However, with material offset in practice, such as chemical composition, strain, interface states being susceptible to such and also the well depth; and quantum size limits of scaling downward in current semiconductor development, where one can reach extremely small physical thickness between layers, where scattering, disorder and phonons amplifies quantization breakage, and forcing experimental conditions to be high (for example, high quality materials). The most important distinction or rather potential point of confusion, justifiably so, is the conjunction between classical, semi-classical dynamics (in which effective mass on mass tensor is an artefact of), and quantum interpretation intersects of material structures and models. Essentially, we are to admit inconsistency and approximated results, with potential realization that epistemically and ontologically, those models are bundles of non-exact interpretation and patterns, thus requires explicit framing of effectiveness, if not to be unable to trace back on failing points or divergent results when applied of different settings. 


\subsection{Quantum potential topology justification}

So far, we have been discussing the potential intervention from external conditioning, leading to boundary condition and wavefunction energy landscape, and so forth. However, one would also add. It is really relevance, or at least that important in determination of the system? Why should one care, if for example, certain conditioning of the theoretical framework pertains to the idea of symmetry and invariance, such is to say of different potential, it is the same physics, thence implicitly allow for potential injection in the form of solving dynamical equation for $V_{1}(q)$, $V_{3}(q)$, enables solving for $V_{1}(q)+V_{2}(q)$? 

Indeed, such change can happen intuitively, and supposedly the general notion of \textbf{topology} suggests a form of continual buildup. That is, given two potential, the topological landscape should operate, in hindsight, similar to a convolution over the arbitrary potential field, leaving us with feasible linear addition on potential landscape. However, such only works if the potential itself is defined as scalar on the configuration space of the dynamical, be it Schrödinger or Dirac, or von Neumann equation, or Heisenberg picture. They change encoding of the supposed system under configuration, but not the system below it. This is the trade-off of quantum mechanics, and one simple cannot ignore this. 

One example that would make sense, is in the discrepancy between quantum mechanical treatment and classical treatment where the philosophy of specification is different. Ontology in classical physics is absolute; one defines properties of the system, and assigns them as given to objects, following by rules under assumptions of interactions, and so on. As said, we say classical physics is externally enforced, and thus bear the limitation of specification frame. That is also why the topic of symmetry was developed during the onset of theoretical mechanics, as a mechanism toward controlling and confining runaway argument of this classical view, while effective, introduces frame relativity, and framework mismatch artefact. Quantum mechanics, on the other side, notice categorical error, on the frame $Q[P]$ upon the physical world $P$ to be insufficient, and often, misaligned to the physical realization itself. Such prompts the reversal action, where one defines $\Psi$ as the placeholder of all sufficient information given of the system, and must be extracted from interaction, or inference, of which is to say the underlying physics is hidden under a black-box, and can only be extracted. This, trades transparency and the capacity of specification on details, toward probabilistic framing and knowledge objectivity, thus allowing quantum physics to take on physical behaviours much more freely and accurately, because of frame shift. Under the same line of argument, we can also understand why Bell's inequality works against hidden variable conjecture. For the hidden variable conjecture to work, it must assume there exists a frame $Q'\supset Q$ previously, that will explain $Q$. However, this brings about the presumption that $Q'$ is the true state of the world, thus before any success can be taken down, $Q'$ simple cannot beat the other side of Bell's inequality, because the inequality itself is forced from the object itself. Any act of external referencing is inconsistent with the physical oracle being the object encoded for decryption (i.e. quantum results), so to speak. This works against locality and realism, where one can say that they are ontological commitments about how information and causation work, and how state with measurement works, respectively, and thus not strictly the physical reality. In together, they assert that "measurement outcomes reference a pre-existing external state of affairs.", which is \textbf{really strong} of an argument \footnote{In essence, what we say is: classical models assume specification completeness. Quantum mechanics assumes interaction-completeness. And finally, the wavefunction is not a state of reality but a sufficient encoding under epistemic constraint, and that is exactly why it works.}. 

That said, such behaviour is typically hard to erase, and there exists a wide-spread problem, where analytical thinking indeed penetrates and assume such stance of the hidden variable interpretation, or rather to fix ontological completeness. For instance, this affects mainly the \textit{analytical method} (or closure of such), if not done correctly, into assuming hidden variable conjecture by itself. This is because of the formation of though where theorems and formal language takes precedence, and thus locking the interpretation in such frame. Because this is a frame, it is treated as similar to saying the formalism is complete as explanation. Which, by itself, invalidates the presumption of quantum mechanics, and Bell's inequality until invoked. Such is also why damage control is often used for justification of models and approximations, as declared and seen in previous section, where one requires justification to apply EMA, and identifying error regions outside EMA's approximation range. Such is also the position on modeling that we will have in mind. 

Moving on, the hidden consequences of quantum physics by then, indicate that the topology itself is rather not scalar, but, the Hamiltonian that encodes such landscape, is an operator. Hence, the potential itself, is expressing a functional space. Such is why while it is intuitive to think of solving $V(x)$ separately, operator group actions, energy discrete levels, transfer topology, and different system depictions change from first order fixture of the Hamiltonian. Thence, unless each potential remains structurally continuous on functional space, which in certain cases are impossible, especially since they are not spectral composition, we cannot solve them separately without introducing additional scaffolding. This is also a problem, and hedges a question of certain modified Hamiltonian that allows spectral decomposition on linearity requirement of the potential topology. Those are conditions of quantum physics. Thence, while acknowledging it as an effective theory, we work within its set of solution consequence, and would reserver such argument on spectral decomposition in another discourse of modified Hamiltonian and the structural presumption needed for an upgraded configuration space. Notably, during $n$th order configuration space addition, one inherits avoidance of solution pathologies, and structural debt.


